\section{Background and Literature Review}
LLMs are increasingly adapted using sensitive domain-specific data in social,
biomedical,
educational,
and research settings.
Differential privacy (DP) provides a rigorous framework for limiting individual-level privacy leakage through calibrated randomized mechanisms \citep{dwork2006calibrating, dwork2014algorithmic}.
In LLM adaptation,
DP can reduce memorization and leakage,
but it also affects optimization,
generation behavior,
and the semantic fidelity of learned or sanitized text.
Much of the private learning literature evaluates this tradeoff primarily through downstream task accuracy or utility,
a summary that leaves alignment degradation unmeasured.

Three lines of work bear on that gap.
Text sanitization and metric DP show that sensitive information can be perturbed while partially preserving semantic structure \citep{chatzikokolakis2013broadening, carvalho2021tem, yue2021differential}.
Research on hallucination and semantic uncertainty argues that reliability failures in LLMs are often better understood at the level of meaning rather than surface token sequences \citep{farquhar2024detecting}.
Semantic privacy research further suggests that privacy leakage can be contextual and inferential,
not limited to verbatim memorization \citep{ma2025sok}.
Each of them locates its object -- the guarantee,
the failure,
the leak -- at the level of meaning rather than of tokens,
which is where an account of what privacy costs has to look as well.

Two features of this literature shape the design of our study.
First,
the mechanisms that could likely be deployed for private adaptation do not share a privacy scale.
DP-SGD attaches an example-level $(\varepsilon, \delta)$ guarantee to the optimization procedure \citep{abadi2016deep, mironov2017renyi},
whereas embedding-space sanitization attaches a word-level metric-LDP guarantee to the text itself \citep{chatzikokolakis2013broadening, feyisetan2020privacy, yue2021differential}.
A budget of,
say,
$\varepsilon = 8$,
therefore means something different in each family,
and a comparison of alignment degradation ``at matched $\varepsilon$'' across families is not defined well.
Any cross-mechanism account of the cost of privacy consequently needs a second axis,
one that is measurable on a common scale and independent of the guarantee that produced it.
Second,
the two families act on different things.
Sanitization acts on meaning by definition,
replacing words with embedding-space neighbors;
DP-SGD leaves the text as-is and instead alters the gradients,
through per-example clipping that biases the update direction plus noise injected per trainable parameter.
Its privacy benefit operates through a different channel,
suppressing the memorization of rare sequences that extraction attacks target \citep{carlini2021extracting},
which is precisely where ``unusual'' facts live.
The families therefore differ not only in what they guarantee but in what they damage,
which is what makes a single privacy axis the wrong instrument for comparing them.
